% ============================================================
%  paquetes.tex — Paquetes utilizados en la tesis
%  Compilación: XeLaTeX + biber (citas IEEE)
% ============================================================

% ----- Idioma -----
\usepackage[spanish,es-tabla]{babel}
\usepackage{csquotes}              % comillas tipográficas (requerido por biblatex)

% ----- Fuentes (XeLaTeX / fontspec) -----
% TeX Gyre Termes es un clon de Times, estándar en TeX Live y MiKTeX.
% Para usar "Times New Roman" sustituye por: \setmainfont{Times New Roman}
\usepackage{fontspec}
\setmainfont{TeX Gyre Termes}      % texto (Times)
\setsansfont{TeX Gyre Heros}       % sans
\setmonofont{TeX Gyre Cursor}      % monoespaciada

% ----- Matemáticas -----
\usepackage{amsmath}
\usepackage{amssymb}
\usepackage{mathtools}
\usepackage{unicode-math}
\setmathfont{TeX Gyre Termes Math} % fuentes de matemáticas acordes a Termes

% ----- Figuras, tablas, unidades -----
\usepackage{graphicx}
\usepackage{float}
\usepackage{longtable}
\usepackage{booktabs}
\usepackage{multirow}
\usepackage{array}
\usepackage{caption}
\usepackage{subcaption}

% ----- Diseño de página -----
\usepackage[a4paper,top=3cm,bottom=3cm,left=3.5cm,right=2.5cm,headheight=14pt]{geometry}
\usepackage{setspace}              % interlineado
\usepackage{fancyhdr}              % encabezados / pies
\usepackage{enumitem}              % listas personalizables
\usepackage{etoolbox}              % adaptación de entradas del índice

% ----- Bibliografía (IEEE + biber) -----
% sorting=none → numeración por orden de cita (típico de IEEE)
\usepackage[style=ieee,backend=biber,sorting=none]{biblatex}

% Adaptar book antes de que hyperref envuelva el comando de capítulo.
% Tabla de contenido con una cabecera por capítulo y páginas automáticas.
\addto\captionsspanish{\renewcommand{\contentsname}{TABLA DE CONTENIDO}}

\makeatletter
\newcommand{\l@chaptergroup}[2]{%
  \addvspace{1em}%
  \@dottedtocline{0}{0em}{0em}{\bfseries #1}{\bfseries #2}}
\newcommand{\toclevel@chaptergroup}{0}
\renewcommand{\l@chapter}[2]{%
  \@dottedtocline{0}{0em}{1.8em}{\bfseries #1}{\bfseries #2}}

% Conserva la numeración normal de book y sus enlaces; solo añade la
% cabecera CAPÍTULO N. antes del título numerado en la tabla de contenido.
\patchcmd{\@chapter}
  {\addcontentsline{toc}{chapter}{\protect\numberline{\thechapter}#1}}
  {\addcontentsline{toc}{chaptergroup}{\protect\MakeUppercase{\@chapapp}\ \thechapter.}%
   \addcontentsline{toc}{chapter}{\protect\numberline{\thechapter}#1}}
  {}
  {\PackageError{tesis}{No se pudo configurar la tabla de contenido}%
    {Comprueba la definicion de chapter en la clase book.}}
\makeatother


% ----- Enlaces / metadatos del PDF (debe ir al final) -----
\usepackage[colorlinks=true,linkcolor=blue,citecolor=blue,urlcolor=blue]{hyperref}
